\section{Kriptanalisis dan Keamanan}
\label{sec:kriptanalisis-knapsack}

\subsection{Ukuran Kunci}

Ukuran cipherteks yang dihasilkan lebih besar daripada plainteksnya. Hal ini
terjadi karena enkripsi dapat menghasilkan kriptogram yang nilai desimalnya
lebih besar daripada nilai desimal blok plainteks yang dienkripsikan.

Untuk menambah kekuatan algoritma knapsack, sumber menganjurkan kunci publik
maupun kunci privat berisi paling sedikit 250 elemen. Setiap elemen sebaiknya
berukuran 200 sampai 400 bit. Sumber juga menyebut modulus berukuran 100 sampai
200 bit.

% Koreksi: anjuran modulus 100--200 bit bertentangan dengan syarat m > jumlah
% seluruh elemen. Untuk 250 elemen yang masing-masingnya paling sedikit 200 bit,
% jumlahnya sudah paling sedikit 449 bit (dihitung dengan Python:
% w1 = 2^199, w_i = jumlah sebelumnya + 1, total.bit_length() = 449).
% Modulus yang benar harus lebih panjang daripada jumlah itu, jadi paling
% sedikit 449 bit, bukan 100--200 bit.
Angka modulus itu tidak konsisten. Pada barisan \textit{superincreasing},
jumlah 250 elemen yang masing-masingnya paling sedikit 200 bit sudah melebihi
$2^{448}$. Modulus harus lebih besar daripada jumlah seluruh elemen, sehingga
panjangnya paling sedikit 449 bit. Modulus 100 sampai 200 bit tidak mungkin
memenuhi syarat tersebut.

\subsection{Serangan Shamir}

Algoritma knapsack dinyatakan sudah tidak aman karena dapat dipecahkan oleh
Adi Shamir. Pada 1982, Shamir merumuskan sebuah transformasi yang memungkinkan
rekonstruksi \textit{superincreasing knapsack} dari \textit{normal knapsack}
dalam waktu polinomial. Makalahnya berjudul \textit{A polynomial time algorithm
for breaking the basic Merkle-Hellman cryptosystem}.

Shamir tidak memecahkan persoalan knapsack secara umum. Ia memanfaatkan struktur
khusus kunci publik Merkle--Hellman. Kunci publik itu merupakan kelipatan
modulo dari sebuah barisan \textit{superincreasing}. Struktur tersebut cukup
untuk memulihkan kunci privat tanpa menelusuri seluruh kemungkinan.
